\clearpage\begingroup\HistoricalHead{229}{R}{BESTIMMUNG DER BAHNEN DER HIMMELSKÖRPER.}
\begin{center}
[8.]\par\medskip
\textls[100]{Dritte Verbesserungsmethode.}
\end{center}
\noindent{\small [S. 7]}\par
\noindent\hspace*{\parindent}Nach drei Hypothesen für Knoten und Neigung berechne man die heliocentrischen Örter für alle drei Beobachtungen und aus diesen den Kegelschnitt. Hiezu wird man folgende Vorschriften vorzüglich brauchbar finden[\textsuperscript{**})]
\begingroup\renewcommand{\thefootnote}{}\footnotetext{[\textsuperscript{**}) Vgl. Theoria motus, art. 82—83, Werke VII, 1906, S. 105.]}\endgroup
{\small
\[
\operatorname{tang}\omega=
\frac{\dfrac{\sin\tfrac12(v''-v')\sin\tfrac12(v''+v')}{r}
+\dfrac{\sin\tfrac12(v'-v)\sin\tfrac12(v'+v)}{r''}
-\dfrac{\sin\tfrac12(v''-v)\sin\tfrac12(v''+v)}{r'}}
{\dfrac{\sin\tfrac12(v''-v')\cos\tfrac12(v''+v')}{r}
+\dfrac{\sin\tfrac12(v'-v)\cos\tfrac12(v'+v)}{r''}
-\dfrac{\sin\tfrac12(v''-v)\cos\tfrac12(v''+v)}{r'}}.
\]
}
\noindent Der Zähler in diesem Ausdruck ist
\[
=2\frac ep\sin\tfrac12(v''-v')\sin\tfrac12(v'-v)\sin\tfrac12(v''-v)\sin\omega,
\]
\noindent der Nenner
\[
=2\frac ep\sin\tfrac12(v''-v')\sin\tfrac12(v'-v)\sin\tfrac12(v''-v)\cos\omega.
\]
\clearpage
\HistoricalHead{230}{L}{THEORETISCHE ASTRONOMIE. NACHLASS.}\SourceLeaf{235}
\noindent\hspace*{\parindent}Man findet also auf diese Weise zu gleicher Zeit $\omega$ und $\dfrac ep$. — Bequemer als obige allgemeine Formel ist folgende, die man leicht aus jener ableitet:
\[
\frac ep\sin\bigl(\omega-\tfrac12(v+v'')\bigr)
=\frac{\dfrac1r-\dfrac1{r''}}{2\sin\tfrac12(v''-v)},
\]
\[
\frac ep\cos\bigl(\omega-\tfrac12(v+v'')\bigr)
=\frac{\begin{aligned}
&\left(\frac1r-\frac1{r'}\right)\operatorname{cotang}\frac12(v'-v)\\
&-\left(\frac1{r'}-\frac1{r''}\right)\operatorname{cotang}\frac12(v''-v')
\end{aligned}}{2\sin\tfrac12(v''-v)}.
\]

Durch Vertauschung des Winkels $v'$ mit $v$ oder $v''$ erhält man aus dieser Formel noch zwei ganz ähnliche. Endlich dient zur Berechnung von $p$ folgende Formel:
\[
\frac1p=\frac{\dfrac1r\sin(v''-v')+\dfrac1{r''}\sin(v'-v)-\dfrac1{r'}\sin(v''-v)}
{4\sin\tfrac12(v''-v')\sin\tfrac12(v'-v)\sin\tfrac12(v''-v)},
\]
\noindent woraus $p$ unmittelbar folgt, oder eine von diesen:
\[
\begin{aligned}
\frac1p&=\frac1r+\frac ep\cos(v-\omega)\\
&=\frac1{r'}+\frac ep\cos(v'-w)\\
&=\frac1{r''}+\frac ep\cos(v''-\omega).
\end{aligned}
\]
\begin{center}[9.]\end{center}
\noindent{\small [S. 8]}\par
\noindent\hspace*{\parindent}Wir fügen diesem Traktat noch einige andere Formeln bei, die bei diesen Rechnungen sehr brauchbar sind:
\[
\begin{aligned}
\sin(v-\lambda)&=\frac{\sin\tfrac12i^2\sin2(v-\AstNode)}{\cos\beta}
=\operatorname{tang}\tfrac12i\sin\beta\cos(\lambda-\AstNode)\\
&=2\sin\tfrac12i^2\sin(v-\AstNode)\cos(\lambda-\AstNode)\\
&=\operatorname{tg}\tfrac12i\operatorname{tg}\beta\cos(v-\AstNode)\\
&=\cos\beta\sin2(\lambda-\AstNode)\mathbin{.}\tfrac12\operatorname{tang}i\mathbin{.}\operatorname{tang}\tfrac12i.
\end{aligned}
\]
\noindent Für die grösste Mittelpunktsgleichung ist $\cos E=-\dfrac{\operatorname{tg}\tfrac12\varphi}{1+\surd\cos\varphi}$ und die
\[
\mathrm{M.\ G.}=2\operatorname{Arc}[\sin]\frac{\sin\tfrac12\varphi\sin E}{\sqrt[4]{\cos\varphi}}+\sin\varphi\sin E.
\]
\noindent Als äusserst scharfe Annäherung ist folgende Formel völlig hinreichend:
\[
\mathrm{Aequ.\ C.\ max.}=\varphi+e+\tfrac38(\varphi-e).
\]
\clearpage
\HistoricalHead{231}{R}{BESTIMMUNG DER BAHNEN DER HIMMELSKÖRPER.}\SourceLeaf{236}
\begin{center}\textls[100]{Differentialformeln:}\end{center}
\noindent\begin{tabular}{@{}l@{\quad}p{0.84\linewidth}@{}}
$m$ . . . . . & mittl. Anomalie zur Zeit der Epoche,\\
$U$ . . . . . & mittl. Bewegung der Sonne von der Epoche bis zur Beobachtung.
\end{tabular}
\[
\begin{aligned}
\mathrm{I.}\quad dv&=d\omega+\frac{aa\cos\varphi}{rr}dm-\frac{3U\cos\varphi}{2rr\surd a}da-\frac ar(2-e\cos M)\sin E\,d\varphi,\\[1ex]
\text{Pars ultima}&=-\frac{aa}{rr}(1+\cos\varphi^2+e\cos E)\sin E\,d\varphi,
\end{aligned}
\]
\[
\mathrm{II.}\quad dr=\left(\frac ra+\frac32\frac{U\operatorname{tg}\varphi\sin M}{a^{\frac32}}\right)da-a\operatorname{tg}\varphi\sin M\,dm+a\cos\varphi\cos M\,d\varphi,
\]
\[
\mathrm{III.}\quad d\lambda=d\AstNode+\frac{\cos i}{\cos\beta^2}d(v-\AstNode)-\cos(\lambda-\AstNode)\operatorname{tg}\beta\,di,
\]
\[
\mathrm{IV.}\quad d\vartheta=\frac{\cos(\lambda-\AstNode)\sin i}{\cos\beta^2}d(v-\AstNode)+\frac{\sin(\lambda-\AstNode)}{\cos\beta^2}di,
\]
\[
\mathrm{V.}\quad d\rho=\cos\beta\,dr-\tfrac12\rho\sin2\beta\,d\vartheta.
\]
\begin{center}\rule{0.25\linewidth}{0.35pt}\end{center}
\clearpage\endgroup
